\documentclass[11pt]{article}
\usepackage[margin=1in]{geometry}
\usepackage{amsmath,amssymb,amsthm}
\usepackage{hyperref}
\newtheorem{theorem}{Theorem}
\newtheorem{lemma}[theorem]{Lemma}
\theoremstyle{definition}
\newtheorem{definition}[theorem]{Definition}
\theoremstyle{remark}
\newtheorem{remark}[theorem]{Remark}
\title{Disproof of Conjecture 00000002160:\\ the Higman--Sims graph does not attain the Hoffman bound}
\author{Submission \texttt{aristotle\_submission\_20261020150000}}
\date{2026-10-20}
\begin{document}
\maketitle

\section{The conjecture}
Conjecture 00000002160 reads:
\begin{quote}
\emph{Definition: The independence number of the Higman--Sims graph. Conjecture: The
independence number equals the Hoffman bound, and the kernel of the bound is the Delsarte
semilattice.}
\end{quote}
The statement is a conjunction. We refute its first, precise, conjunct; hence the whole
conjunction is false. (The second conjunct, ``the kernel of the bound is the Delsarte
semilattice'', has no standard mathematical meaning and is not used.)

\section{Definitions}
\begin{definition}[Hoffman bound]
Let $G$ be a $k$-regular graph on $n$ vertices whose adjacency matrix $A$ has least eigenvalue
$\lambda_{\min}$. Hoffman's ratio bound is
\[ H(G) \;=\; \frac{n\,(-\lambda_{\min})}{k-\lambda_{\min}} . \]
Hoffman's theorem states $\alpha(G)\le H(G)$, where $\alpha$ is the independence number;
``$\alpha(G)$ equals the Hoffman bound'' (equivalently, ``the Hoffman bound is attained'')
means $\alpha(G)=H(G)$.
\end{definition}

\begin{definition}[Higman--Sims graph]
Let $\mathcal{S}$ be the Steiner system $S(3,6,22)$: $22$ points and $77$ six-element blocks
(hexads) such that any $3$ points lie in exactly one hexad. The Higman--Sims graph $HS$ has
the $100$ vertices $\{\infty\}\cup\{\text{points}\}\cup\{\text{hexads}\}$; $\infty$ is adjacent to
every point, a point is adjacent to the hexads containing it, and two hexads are adjacent iff
they are disjoint. It is the (unique) strongly regular graph with parameters
$(100,22,0,6)$.
\end{definition}

\section{The disproof}
\begin{theorem}\label{thm:main}
$HS$ is $22$-regular, its least adjacency eigenvalue is $-8$, and
$H(HS) = \frac{100\cdot 8}{22+8} = \frac{80}{3}$. Consequently $\alpha(HS)\neq H(HS)$, i.e.
the conjecture is false.
\end{theorem}

\begin{lemma}\label{lem:eig}
Let $G$ be any strongly regular graph with parameters $(100,22,0,6)$ and adjacency matrix $A$.
Then every eigenvalue of $A$ lies in $\{22,2,-8\}$, and $-8$ is an eigenvalue.
\end{lemma}
\begin{proof}
Strong regularity says $A^2 = 22I + 0\cdot A + 6(J-I-A)$, where $J$ is the all-ones matrix.
Let $Av=\mu v$ with $v\neq 0$ and put $s=\sum_x v_x$. Applying the identity to $v$ gives
$(\mu^2+6\mu-16)\,v = 6s\,\mathbf 1$. If $\mu^2+6\mu-16\ne 0$, then $v$ is a nonzero constant
vector, and since $G$ is $22$-regular, $A v = 22 v$, so $\mu=22$. Otherwise
$(\mu-2)(\mu+8)=0$. Hence every eigenvalue is $22$, $2$ or $-8$.

The eigenvalues $\mu_1,\dots,\mu_{100}$ (with multiplicity) sum to $\operatorname{tr}A=0$. If none
of them were $-8$, each would be $\ge 2$ and their sum would be $\ge 200>0$, a contradiction.
\end{proof}

\begin{proof}[Proof of Theorem~\ref{thm:main}]
By Lemma~\ref{lem:eig}, $\lambda_{\min}=-8$, and $k=22$, $n=100$, so $H(HS)=800/30=80/3$.
The independence number is a natural number, and $3\alpha = 80$ has no integer solution; hence
$\alpha(HS)\ne 80/3$.
\end{proof}

\begin{remark}
The neighbourhood of $\infty$ (the $22$ points) is independent, so $\alpha(HS)\ge 22$; this is
also checked in Lean. It is known that $\alpha(HS)=22$; this fact is \emph{not} needed and not
formalised here. Under the alternative reading ``$\alpha(HS)=\lfloor H(HS)\rfloor=26$'' the
conjecture would also be false by this known value, but that reading is not covered by the
formal proof.
\end{remark}

\section{Lean formalisation}
The Lean~4 project in \texttt{lean/} (Lean 4.28.0, Mathlib \texttt{v4.28.0}) contains:
\begin{itemize}
\item \texttt{Conjecture2160/HigmanSims.lean}: the $77$ hexads of $S(3,6,22)$ as $22$-bit masks
(derived from the extended binary Golay code), a kernel-checked proof
(\texttt{hexads\_steiner}) that they form an $S(3,6,22)$, the graph
\texttt{higmanSims : SimpleGraph (Fin 100)} defined exactly as above (vertex $0=\infty$,
$1,\dots,22$ = points, $23,\dots,99$ = hexads), and \texttt{higmanSims\_isSRG}:
\texttt{higmanSims.IsSRGWith 100 22 0 6} (Mathlib's notion). The latter is proved by
\texttt{decide +kernel} on a Boolean checker together with a soundness proof linking the
checker to Mathlib's \texttt{commonNeighbors} and \texttt{degree}.
\item \texttt{Conjecture2160/Spectral.lean}: \texttt{minAdjEigenvalue} (least eigenvalue of the
real symmetric adjacency matrix, via Mathlib's spectral theorem), \texttt{hoffmanBound}, and
Lemma~\ref{lem:eig} for every srg$(100,22,0,6)$
(\texttt{srg\_minAdjEigenvalue}, \texttt{srg\_hoffmanBound}, \texttt{srg\_indepNum\_ne\_hoffmanBound}).
\item \texttt{Conjecture2160/Main.lean}: the main theorem
\texttt{conjecture\_2160\_false}, stating
\texttt{higmanSims.IsRegularOfDegree 22} and
\((\texttt{higmanSims.indepNum} : \mathbb{R}) \neq \texttt{hoffmanBound higmanSims 22}\),
and \texttt{higmanSims\_hoffmanBound : hoffmanBound higmanSims 22 = 80/3}.
\end{itemize}
All theorems depend only on the axioms \texttt{propext}, \texttt{Classical.choice},
\texttt{Quot.sound}; there is no \texttt{sorry} and no \texttt{native\_decide}. The auxiliary
script \texttt{gen\_hs.py} regenerates the hexad and row data and independently re-checks the
Steiner and strong-regularity properties.

\end{document}
